% TotalSegmentator -> start points -> in-house tracer (Kit Mills Bransby). Not the author's
% contribution, so only outputs and how it was run on CGPS; the author's check is sec:method-medis.
% Questions for Kit are in the red note below; write the section from his answers, then delete the note.
% Red notes need \usepackage{xcolor}. Remove before submission.

\newpage

\providecommand{\clarify}[1]{\textcolor{red}{\textbf{Clarification needed:} #1}}
\providecommand{\ask}[1]{\textcolor{red}{[#1]}}
\section{Coronary Centerline and Lumen Extraction}
\label{sec:method-tracer}

To extract the coronary centerlines and lumen from the baseline CCTA volumes, an in-house tracer was applied to all scans of the Herlev--{\O}sterbro cohort (Section~\ref{sec:data-cohort}). The in-house tracer was developed and run by Kit Mills Bransby. The method is described here so that the full pipeline is shown. 


\clarify{Questions for Kit (for a description of 1--2 pages).\\
\textit{Method}\\
-- Overall pipeline, from CT volume to centerlines and lumen surface, in a few steps.\\
-- Name of the method: what should I call it instead of ``in-house tracer''?\\
-- Type of network (e.g.\ nnU-Net, 3D U-Net), its size and main settings.\\
\textit{Training}\\
-- Exactly which data (datasets, number of scans, scanners)?\\
-- How were the labels made, and what are they (only artery vs background, or more classes)?\\
\textit{Output}\\
-- Are the centerlines made from a predicted segmentation, or predicted directly? The final centerlines will have smoothing right? How is that made?\\
-- How are the arteries labelled (LM, LAD, LCx and so on)?\\
-- How is the lumen surface made?\\
\textit{Performance}\\
-- Did you measure the performance, so I have an upper bound of how good we expect the predictions to be}